Trong môi trường \emph{multi-tenant Kubernetes}, nhiều người thuê (tenant) bị cô
lập bằng không gian tên (namespace) nhưng vẫn \emph{chia sẻ chung một nhân} hệ điều
hành; do đó lời gọi hệ thống (syscall) trở thành bề mặt tấn công trọng yếu cho
\emph{leo thang đặc quyền} và thoát container---kẻ tấn công vượt qua ranh giới
tenant để chiếm quyền trên máy chủ. Báo cáo chuyên đề này nghiên cứu một
\emph{cơ chế phát hiện} các tấn công leo thang đặc quyền dựa trên phân tích syscall
thời gian thực. Cơ chế kế thừa thiết kế lõi phát hiện bất thường của công trình
tham chiếu DeSFAM~\cite{desfam_original} (SyscallAD: Bộ mã hóa biến phân tự động VAE
hợp nhất với Rừng cô lập iForest), trong đó các đặc trưng \emph{hiện diện của syscall
đặc quyền} (\texttt{ptrace}, \texttt{setuid/capset}, \texttt{setns},
\texttt{unshare}, \texttt{mount}, \texttt{bpf}, nạp mô-đun nhân) chính là tín hiệu
hành vi của leo thang đặc quyền. Mô hình được huấn luyện trên bộ dữ liệu công khai
DongTing~\cite{ref26_dongtingdataset} và được hiện thực hóa thành một bộ phát hiện
chạy trực tiếp: cảm biến syscall dùng Tetragon~\cite{tetragon} (eBPF) giám sát theo
namespace trên cụm Kubernetes, đẩy cảnh báo qua Prometheus/Grafana. Toàn bộ quy
trình huấn luyện và suy luận được đóng gói bằng Docker để bảo đảm khả năng tái lập.
Trên DongTing, mô hình VAE đạt AUC kiểm định $0{,}944$; tập hợp VAE\,+\,iForest đạt
AUC kiểm định $\approx 0{,}95$ và AUC trên tập kiểm tra $0{,}835$ (F1 $0{,}417$ tại
ngưỡng bảo thủ chọn trên tập kiểm định). Khi triển khai trực tiếp, điểm bất thường
tách biệt rõ giữa khối lượng công việc lành tính ($\le 0{,}19$) và các kịch bản leo
thang đặc quyền/thoát container mô phỏng theo các họ CVE thực
(\mbox{CVE-2021-4034}, \mbox{CVE-2022-0847}, \mbox{CVE-2019-5736}) trong namespace
thử nghiệm ($0{,}278$--$0{,}930$), cho phép một ngưỡng vận hành $T_0 = 0{,}25$ phát
hiện tấn công thời gian thực. Báo cáo cũng đối chiếu trung thực cơ chế đề xuất với
số liệu công bố của DeSFAM, chỉ ra các điểm tương đồng và khác biệt, cùng bài học
then chốt về căn chỉnh không gian đặc trưng giữa huấn luyện và phục vụ.
